\documentclass[11pt]{article}
\usepackage{fontspec}
\usepackage[french]{babel}
\usepackage{amsmath}
\usepackage{amssymb}
\usepackage{amsthm}
\usepackage{mathrsfs}
\usepackage{tikz-cd}
\usepackage{enumitem}

\newcommand{\printedpage}[1]{\marginpar{\footnotesize p.~#1}}
\newcommand{\msp}[1]{\marginpar{\footnotesize [#1]}}
% printed notation
\newcommand{\M}{\mathcal{M}}
\newcommand{\Cat}{\mathrm{Cat}}
\newcommand{\Ho}{\mathrm{Ho}}
\newcommand{\Fl}{\mathrm{Fl}}
\newcommand{\Cof}{\mathrm{Cof}}
\newcommand{\Diag}{\mathrm{Diag}}
\newcommand{\id}{\mathrm{id}}
\newcommand{\deq}{\overset{\text{déf}}{=}}
% small tilde printed after a symbol (Q-saturation)
\newcommand{\Qs}{\tilde{\ }}
% arrow with one head on the left and a double head on the right (factorisation)
\newcommand{\fact}{\mathrel{\leftarrow\mkern-11mu\twoheadrightarrow}}
% arrow with half heads (harpoons) at both ends, as printed on p. 23
\newcommand{\harp}{\mathrel{\leftharpoondown\joinrel\relbar\joinrel\rightharpoonup}}

\begin{document}

\printedpage{19}
% head: A. Grothendieck : Dérivateurs, Ch. XIX, Édition des Universités de Montpellier II et Paris VII
% foot: Version 11 février 2019

\msp{page 31}
\textbf{Scholie 2.10.} \textit{Soit $(W,C,F)$ un triple de Quillen clos dans $\M$. Pour $I$ dans \textup{Cat}, pour que $(W(I),C\langle I\rangle,F(I))$ soit un triple de Quillen clos pour $\M(I)$, il faut et il suffit que les couples $C\langle I\rangle\longleftrightarrow TF(I)$ et $TC\langle I\rangle\longleftrightarrow F(I)$ soient factorisants. Supposons que cette condition soit satisfaite pour $I$ et pour $I'$, et soit $u : I\longrightarrow I'$ une flèche de $I$ dans $I'$. Alors les assertions de \textup{2.3. 2°} sont valables, en particulier on trouve un foncteur $u^{W}_{!} : \Ho_W(I)\longrightarrow\Ho_W(I')$, adjoint à gauche de $u^{*}_{W}$\footnote[18]{Dire aussi que le calcul des fibres de $u^{W}_{!}(X)$, prévu dans Dér 5 bis, est OK.}. Ces résultats se dualisent aussitôt (par $\M\longmapsto\M^{\circ}$, $I\longmapsto I^{\circ}$) pour les triples $(W(I),C(I),F[I])$, et pour la construction de $u^{W}_{*}$, adjoint à droite de $u^{*}_{W}$. En particulier, pour que le triple précédent soit un triple de Quillen (forcément clos), il faut et il suffit que les deux couples $C(I)\longleftrightarrow TF[I]$ et $TC(I)\longleftrightarrow F[I]$ soient factorisants ; et quand cette condition est satisfaite pour $I$ et pour $I'$ (pour une flèche $u : I\longrightarrow I'$ dans \textup{Cat} donnée), alors $u^{W}_{*}$ existe et se construit dualement de la construction de $u^{W}_{!}$ explicitée dans \textup{2.3. 2°}.}

\msp{page 32}
% a vertical bar is printed in the left margin alongside Corollaire 2.11
\textbf{Corollaire 2.11.} \textit{Le prédérivateur $\mathbf{D}_{\M,W}$ défini pour un localiseur \textup{très spécial} $W$ dans $\M$ est un dérivateur en plein sens du terme, i.e. il satisfait à tous les axiomes Dér 1 à Dér 6 (y inclus Dér 4 bis et Dér 5 bis et Dér 6 bis) des dérivateurs.}

\textbf{N.B.} À ma propre surprise, je n'ai pas été capable de trouver un cas d'un localiseur $W$ qui définisse un dérivateur sur Cat tout entier (avec existence de $u^{W}_{!}$, $u^{W}_{*}$), \textit{que} pour $W$ un localiseur très spécial, donc s'insérant dans un triple de Quillen très spécial -- ou quand on peut s'y ramener de la façon standard. (Par exemple dans le cas où $\M=\Cat$, $W$ un localiseur fondamental satisfaisant Loc 6 a) et Loc 8, 8 bis)\footnote[19]{axiome des carrés cocartésiens (version directe), axiome des limites inductives, accessibilité.} -- les monomorphismes dans $W$ se refusent formellement à être stable par cochangement de base !) J'ignore totalement si le résultat mirobolant ``prouvé'' par \textsc{Anderson} est vrai ou faux. Ça a l'air presque trop beau pour être vrai \ldots{}D'ailleurs, je ne connais pas de structure de Quillen spéciale qui ne soit très spéciale -- désolé ! % [sic] stable -> stables, agreement with « les monomorphismes »

\msp{page 33}

\section*{3\quad Cas des catégories à cofibrations, des catégories à fibrations}

Je rappelle les axiomes d'une \textit{structure à cofibrations} $(W,C)$ sur une catégorie $\M$.

\begin{itemize}[leftmargin=4em,labelwidth=3.5em,align=left]
\item[Cof$_0$ :] $W$ est un localiseur.
\item[Cof$_1$ :] $C$ contient les isomorphismes, est stable par cochangement de base, par composition.
\item[Cof$_2$ :] $C\cap W$ est stable par cochangement de base (et aussi par composition, et contient les isomorphismes, en vertu de Cof$_0$, Cof$_1$).
\item[Cof$_3$ :] Toute flèche $f$ dans $\M$ se factorise en $f=pi$, avec $i\in C$, $p\in W$.
\end{itemize}

(Axiomes de K. S. \textsc{Brown}, légèrement généralisés par \textsc{Anderson}). Il y a un ``axiome
\printedpage{20}
% head: A. Grothendieck : Dérivateurs, Ch. XIX, Édition des Universités de Montpellier II et Paris VII
% foot: Version 11 février 2019
de continuité'' ou \textit{``axiome transfini''} facultatif, qui est le suivant :

\begin{itemize}[leftmargin=4em,labelwidth=3.5em,align=left]
\item[Cof$_4$ :] $\M$ stable par petites $\varinjlim$ filtrantes, et $C$ et $TC = C\cap W$ sont stables par composition transfinie.
\end{itemize}

Ainsi, Cof$_1$, Cof$_2$, Cof$_4$ peuvent se résumer en disant que $C$ et $TC$ sont Q-saturés transfiniment à gauche.

La structure de cofibration est dite \textit{spéciale} si elle satisfait Cof$_4$, si $\M$ est spéciale, et si $C$ admet une \textit{petite} partie $C_0$ qui soit Q-génératrice à gauche, i.e. telle que
\[
C \ \subset\ C_0\Qs\,, \qquad \text{i.e.}\quad \tilde{C}=C_0\Qs. \tag*{(3.1)}
\]

\msp{page 34}
Alors, par le théorème de factorisation 1.1, on a
\[
C \ \fact\ TF \ \deq\ C_* \ ,
\]
et l'axiome de factorisation Cof$_3$ implique (cf. 1.4) $TF\subset W^{\natural}$. Plus généralement, on a :

\textbf{Proposition 3.1.} \textit{Supposons que $\M$ soit spéciale (stable par petites $\varinjlim$ , et éléments accessibles), que $(W,C)$ satisfasse à Cof$_0$, Cof$_1$, et que $C$ soit stable par composition transfinie (ce qui, joint à Cof$_1$, s'énonce en disant que $C$ est Q-saturé transfini à gauche), enfin que $C$ admette une petite partie $C_0$ Q-génératrice à gauche, de sorte que $C\fact TF\deq C_*$. Alors pour que l'axiome de factorisation Cof$_3$ soit satisfait, il suffit qu'on ait}
\[
TF \ \subset\ W \ , \tag*{(3.1.1)}
\]
\textit{et il faut qu'on ait $TF\subset W^{\natural}$, donc si $W=W^{\natural}$, l'axiome Cof$_3$ équivaut à} (3.1.1).

La structure de cofibration est dite \textit{semi-close} si $C$ est Q-clos à gauche, i.e. si on a $C=C\Qs$ (ce qui implique que $C$ est stable par composition transfinie), \textit{close} si de plus $TC=(TC)\Qs$ (ce qui implique l'axiome transfini Cof$_4$).

\msp{page 35}
Je ne reviens pas ici sur la théorie de K. S. \textsc{Brown} des catégories à cofibrations. C'est à ma connaissance l'ensemble d'axiomes le plus économique connu, pour assurer d'une part une bonne prise sur $\Ho_W$ (par une notion correspondante d'homotopie d'applications, et un calcul de fractions dans $\Pi\M_{\mathrm{c}}$), et d'autre part l'existence de la suite \textit{exacte} de cofibrations, donc une théorie de cofibres $W$-homotopiques. (\textbf{N.B.} Quand la structure est \textit{propre,} i.e. $C\subset\Cof_W$, alors l'existence de cette structure revient à dire simplement que ``$W$ admet assez de $W$-cofibrations'', i.e. que toute $f\in\Fl(\M)$ se factorise en $pi$, avec $i\in\Cof_W$, $p\in W$ ; et on peut donc prendre $C=\Cof_W$. Dans le cas où $W$ est stable par petites $\varinjlim$ filtrantes, $C=\Cof_W$ et $TC=W^{\mathrm{univ}}$ satisfont de plus l'axiome transfini Cof$_4$.) Enfin, K. S. \textsc{Brown} montre que toute flèche $i : X\longrightarrow Y$ dans $C$, \textit{avec $X$, $Y$ cofibrants} (i.e. $\emptyset_{\M}\longrightarrow X$, $\emptyset_{\M}\longrightarrow Y$ dans $C$) est dans $\Cof_W$ (``Lemme de Brown'').

Je m'intéresse ici à la question d'étendre la structure à cofibrations à une catégorie de la forme $\M(I)$, et d'utiliser cette extension pour construire
\msp{page 36}
les foncteurs $u^{W}_{!}$. Notons déjà que la théorie de K. S. \textsc{Brown}, avec les compléments d'\textsc{Anderson}\footnote[20]{revus et arrangés par mes soins.}, donne essentiellement ceci :

\printedpage{21}
% head: A. Grothendieck : Dérivateurs, Ch. XIX, Édition des Universités de Montpellier II et Paris VII
% foot: Version 11 février 2019

\textbf{Proposition 3.2.} \textit{Soit $W\subset\Fl(\M)$ un localiseur associé à une structure à cofibrations $(C,W)$ dans $\M$.}
\begin{itemize}[leftmargin=3em,labelwidth=2.5em,align=right]
\item[1°)] \textit{Sur la catégorie \textup{$\Diag_0$} des ensembles ordonnés \textup{finis}, ce prédérivateur satisfait aux axiomes Dér 1 (décomposition), Dér 2 (localisation), Dér 3 (objet final), Dér 4 bis (existence de $u^{W}_{!}$), Dér 5 bis (calcul des fibres de $u^{W}_{!}$), Dér 6 bis (exactitude à gauche).}
\item[2°)] \textit{C'est encore vrai sur la sous-catégorie \textup{$\Diag_1$} dans \textup{Cat} plus grande, formée des ensembles ordonnés $I$ \textup{noethériens} tels que l'espace topologique associé à $I^{\circ}$ soit noethérien (toute suite croissante d'objets de $I$, ou d'ouverts de $I^{\circ}$, est stationnaire).}
\item[3°)] \textit{Enfin, si l'axiome transfini Cof$_4$ est satisfait, i.e. si $C$ et $TC$ sont stables par composition transfinie, alors la conclusion de \textup{1°} reste valable sur la sous-catégorie pleine \textup{$\Diag_2$} de \textup{Cat} formée des ensembles ordonnés qui sont noethériens sans plus.}
\end{itemize}

\textbf{Corollaire 3.3.} \textit{Supposons que l'on ait Cof$_3$ sous la forme plus forte d'une factorisation
\msp{page 37}
$f=pi$ \textup{fonctorielle} en $f$, ce qui est le cas en particulier si la structure à cofibrations est spéciale. Alors pour toute $I$ dans \textup{Cat}, $\M(I)$ muni de $(W(I),C(I))$ est une catégorie à cofibrations, laquelle satisfait l'axiome transfini Cof$_4$ si $(W,C)$ y satisfait. Donc on peut lui appliquer les résultats de \textup{3.2}, qui impliquent notamment que les axiomes Dér 1, 2, 3 ( ?) sont valables sur \textup{Cat} tout entier, et que $u^{W}_{!}$ existe quand $u$ est de la forme $\id_I\times v$, avec $v : J\longrightarrow K$ une flèche dans \textup{$\Diag_1$} (resp. dans \textup{$\Diag_2$}, si l'axiome transfini est satisfait).}

\textbf{3.4.} La clef de la théorie est dans l'extension de la structure à cofibrations $(W,C)$ de $\M$ à une catégorie $\M(I)$, où $I$ est un ensemble ordonné, en $(W(I),C_I)$. Les faits principaux sont énoncés dans XVII, dans les théorèmes p. 26-27, p. 50, p. 60. C'est quand $I$ est \textit{noethérien,} et qu'on suppose soit que $\mathrm{Esp}(I^{\circ})$ est également noethérien (p.ex. $I$ fini), soit que $(W,C)$ satisfait l'axiome transfini Cof$_4$, qu'on prouve que $(W(I),C_I)$ est également une structure à cofibrations sur $\M(I)$.

\msp{page 38}
Un fait important pour la construction des foncteurs $u^{W}_{!}$, pour $u : I\longrightarrow I'$ une flèche entre catégories ordonnées noethériennes, est la relation
\[
C_I\cap W(I) \ = \ (TC)_I \ , \tag*{(3.4.1)}
\]
d'où résulte que l'on a, pour $u$ comme dessus,
\[
u_!(C_I\cap W(I)) \ \subset\ C_{I'}\cap W(I') \ , \tag*{(3.4.2)}
\]
ingrédient essentiel de la construction du foncteur $u^{W}_{!}$. Pour étendre cette construction au cas d'une flèche $\pm$ quelconque dans Cat, il se pose donc le problème d'étendre la définition de $C_I$, et les formules (3.4.1) et (3.4.2), au cas de catégories $\pm$ quelconques dans Cat. Or on a prouvé :

\textbf{Proposition 3.5.} \textit{Si $\Phi=\Psi^*$ est une partie \textup{Q-close} à gauche de \textup{$\Fl(\M)$}, et $I$ un ensemble ordonné noethérien, $\Phi_I$ l'extension de $\Phi$ en une partie de \textup{$\Fl(\M(I))$} définie par \textup{\textsc{Anderson}}, alors on a}
\[
\Phi_I \ = \ \Psi(I)^* \ = \ v_!(\Phi(I_0))\Qs \ , \tag*{(3.5.1)}
\]
\textit{où $v : I_0\longrightarrow I$ est le foncteur canonique de la catégorie discrète définie par \textup{$\mathrm{Ob}\,I$}.}

\printedpage{22}
% head: A. Grothendieck : Dérivateurs, Ch. XIX, Édition des Universités de Montpellier II et Paris VII
% foot: Version 11 février 2019

C'est (avec des notations différentes) le lemme 6.2.5 (p. 107) de XVII, dont je ne reprends pas ici la démonstration.

Ceci donne la clef d'une définition de $\Phi_I$ comme $\Phi\langle I\rangle$ (avec les notations des paragraphes 1, 2),
\msp{page 39}
pour $I$ quelconque dans Cat,
\[
\Phi_I \ \deq\ \Phi\langle I\rangle \ \deq\ \Psi(I)^* \ = \ v_!(\Phi(I_0))\Qs \ , \tag*{(3.5.2)}
\]
où $\Psi=\Phi_*$, et où on suppose $\Phi$ Q-clos à gauche, i.e. $\Phi=\Psi^*$. Dans le cas où on suppose seulement $\Phi$ Q-saturé transfini sans plus, on aurait envie de poser
\[
\Phi_I \ \deq\ \text{Q-saturé transfini de } v_!(\Phi(I_0)) \ , \tag*{(3.5.3)}
\]
et je présume que pour $I$ ordonné noethérien, cela coïncide avec la définition d'\textsc{Anderson}. Mais je vais me borner pour la suite au cas où $C=\Phi$ est Q-clos à gauche.

Revenons alors au cas d'une structure à cofibrations $(W,C)$ dans $\M$, en supposant $C$ et $TC=C\cap W$ Q-clos à gauche. Définissons pour toute $I$ dans Cat $C_I=C\langle I\rangle$ comme dans (3.5.2). On a alors, pour $u : I\longrightarrow I'$ une flèche dans Cat,
\[
u_!(C\langle I\rangle) \ \subset\ C\langle I'\rangle \ , \tag*{(3.5.4)}
\]
mais on n'y a guère de raison d'espérer avoir aussi (3.4.2), sans disposer de la formule
\[
C\langle I\rangle\cap W(I) \ \overset{?}{=}\ (TC)\langle I\rangle \qquad\qquad \text{\footnote[21]{\textbf{N.B.} On a l'inclusion $\supset$, car $v^*v_!TC(I_0)\subset TC(I_0)$.}}. \tag*{(3.5.5)}
\]
On a prouvé que cette formule est valable si $I$ est ordonné noethérien (XVII th. 2.1, p. 26-27). Cela
\msp{page 40}
nous amène à la

\textbf{Question 3.6.} \textit{Soit $\M$ une catégorie stable par petites limites, et soit $(W,C)$ une structure à cofibrations dans $\M$ telle que $C$ et $TC=C\cap W$ soient Q-clos à gauche. Soit $I$ un objet de \textup{Cat}, a-t-on alors}
\[
C\langle I\rangle\cap W(I) \ = \ (TC)\langle I\rangle \text{ ?} \tag*{(3.6.1)}
\]

On a vu au paragraphe précédent qu'il en est ainsi si $(W,C)$ s'insère dans un triple de Quillen clos $(W,C,F=(TC)_*)$, et quand on suppose de plus que les couples
\[
C\langle I\rangle \ \longleftrightarrow\ TF(I) \ , \qquad TC\langle I\rangle \ \longleftrightarrow\ F(I)
\]
sont factorisant, auquel cas $(W(I),C\langle I\rangle,F(I))$ est une structure de Quillen close dans $\M(I)$. D'ailleurs, $(W,C,F)$ est une structure de Quillen (nécessairement close) si et seulement si les couples $C\longleftrightarrow \underbrace{TF}_{\deq C_*}$, $TC\longleftrightarrow \underbrace{\ F\ }_{\deq (TC)_*}$ eux-mêmes sont factorisants, et si % [sic] factorisant -> factorisants, agreement with « les couples »
\printedpage{23}
% head: A. Grothendieck : Dérivateurs, Ch. XIX, Édition des Universités de Montpellier II et Paris VII
% foot: Version 11 février 2019
de plus $TF\subset W$ (ce qui est automatique si $W=W^{\natural}$, car on aura en tous cas, par Cof$_4$, $TF\subset W^{\natural}$).

Mais peut-on s'en tirer avec moins, pour pouvoir conclure à (3.6.1) ? Ne suffirait-il pas de savoir que $TC\langle I\rangle\longleftrightarrow F(I)$ est factorisant (ce qui n'a guère de chances d'être vrai pour une $I$ $\pm$ générale, si on n'a déjà $TC\fact F$, factorisation dans $\M$ lui-même). C'est bien le cas :

\msp{page 41}
\textbf{Proposition 3.7.} \textit{Sous les conditions générales de \textup{(3.6)} sur $\M$, $W$, $C$, $I$, supposons que le couple $TC\langle I\rangle\longleftrightarrow F(I)$ (où $F\deq (TC)_*$) soit factorisant, i.e. $TC\langle I\rangle\fact F(I)$. Alors la formule \textup{(3.6.1)} est vraie.}

En effet, si $f\in C\langle I\rangle\cap W(I)$, on factorise $f$ en $pi$, avec $i\in TC\langle I\rangle$, $p\in F(I)$. Comme $f,i\in W(I)$, on aura $p\in W(I)$, i.e. $p\in TF(I)$, et à cause de $\underbrace{C\langle I\rangle}_{\ni\, f}\harp\underbrace{TF(I)}_{\ni\, p}$, il existe $p : Y\longrightarrow X'$ telle que $pq=\id_Y$, $qf=i$, % reading resolved: DER_XIX_corrections.csv; % [sic] p : Y -> q : Y, the arrow Y -> X' is q (see the diagram and pq = id_Y)
\[
\begin{tikzcd}[column sep=small, row sep=large]
X \arrow[rr, "f\in C\langle I\rangle\cap W(I)"] \arrow[dr, "i\in TC\langle I\rangle"'] & & Y \arrow[dl, dotted, shift right=1ex, "q"'] \\
& X' \arrow[ur, "p\in F(I)"'] \ , &
\end{tikzcd}
\]
donc $f$ est facteur direct de $i\in TC\langle I\rangle$. Comme $TC\langle I\rangle$ est stable par facteurs directs, on gagne.

\textbf{Corollaire 3.8.} \textit{Soit $(W,C,F)$ une structure de Quillen close dans $\M$, et supposons que $TC=C\cap W$ soit une partie Q-spéciale de $\M$, i.e. le couple $TC\fact F$ un couple spécial (i.e. $\M$ est \textup{spéciale}, et il existe une petite partie $TC_0$ Q-génératrice à gauche pour $TC$). Alors pour toute $I$ dans \textup{Cat}, on a}
\[
C\langle I\rangle\cap W(I) \ = \ TC\langle I\rangle \ . \tag*{(3.8.1)}
\]

\msp{page 42}
En apparence, on n'a pas eu à supposer que $C$ soit elle aussi une partie Q-spéciale de $\Fl(\M)$. Mais comme on a besoin de la factorisation pour $C\longleftrightarrow TF$ (pour avoir au moins une structure de cofibrations sur $\M$), on n'a guère le moyen de l'établir sans supposer $C$ spéciale. Sans compter que pour pouvoir expliciter (3.8.1) pour la construction de $u^{W}_{!}$, on a besoin de savoir que $(W(I),C\langle I\rangle)$ définit bien une structure à cofibrations sur $\M(I)$, ce qui ne peut guère se faire (il me semble) sans avoir la factorisation pour $C\langle I\rangle\longleftrightarrow TF(I)$ également ; chose que je n'arrive à établir pour un $I$ général dans Cat, que moyennant l'hypothèse de spécialité sur $C$ également.

\textbf{3.9.} Pour en avoir le cœur net, il faudrait quelques exemples concrets et illuminants, du type suivant :

a) \texttt{[Il n'y a pas de b).]} Structure de Quillen close $(W,C,F)$ sur $\M$, telle que pour le prédérivateur associé, $u^{W}_{!}$ n'existe pas toujours, ou $u^{W}_{*}$ n'existe par toujours. Je pense notamment au cas d'une flèche $u : I\longrightarrow I'$ avec $I'=e$ (présumant que si dans ce cas, $u^{W}_{!}$ resp. $u^{W}_{*}$ existe, alors il existe dans tous les cas). Un tel exemple serait encore plus % [sic] n'existe par toujours -> n'existe pas toujours, typing error

\end{document}
